\part{Der Zufall im Rechner}

\chapter{Ein Experiment im Rechner}

\section{Der Leitfall}

\begin{merksatz}
\boxtitel{Zwölf Betten}
Eine Intensivstation: Patienten treffen im Mittel $2{,}4$-mal am Tag ein, unabhängig voneinander. Die Verweildauer ist logarithmisch normalverteilt mit Median vier Tage und Streuung $0{,}6$ auf der logarithmischen Skala --- im Mittel also $4{,}79$ Tage. Zwölf Betten.

\begin{center}
\begin{minipage}{0.9\textwidth}
\textbf{Wie oft muss ein Patient abgewiesen werden, weil alle Betten belegt sind?}
\end{minipage}
\end{center}

\schritt{Die Formel}
Das System ist eine Warteschlange ohne Warteraum --- Band 25, Kapitel~5. Für sie gilt die Formel von Erlang, und bemerkenswerterweise hängt sie nur vom Mittel der Verweildauer ab, nicht von ihrer Form:
\[ a = \lambda\cdot\bar d = 2{,}4\cdot4{,}79 = 11{,}49, \qquad \Pr[\text{Abweisung}] = \frac{a^{c}/c!}{\sum_{k=0}^{c} a^{k}/k!} = 0{,}1785 \]

\begin{center}
\begin{minipage}{0.88\textwidth}
\kursiv{\textbf{Fast jeder fünfte Patient findet kein Bett.} Das ist das exakte Ergebnis, und es wird gleich zum Maßstab: Eine Simulation, die diese Zahl nicht trifft, ist falsch. Danach verlassen wir den Sonderfall, in dem es eine Formel gibt.}
\end{minipage}
\end{center}
\end{merksatz}

\begin{merksatz}
\boxtitel{Dasselbe als Simulation}
Man erzeuge Ankünfte, ziehe Verweildauern, führe Buch über die Betten, zähle Abweisungen:

\begin{center}
{\small
\renewcommand{\arraystretch}{1.48}
\begin{tabular}{@{}P{2.0cm}P{2.2cm}P{2.4cm}P{1.8cm}@{}}
\toprule
\textbf{Tage} & \textbf{Ankünfte} & \textbf{Abweisung} & \textbf{Fehler}\\
\midrule
$100$ & $238$ & $0{,}160$ & $0{,}019$\\
$1\,000$ & $2\,387$ & $0{,}174$ & $0{,}005$\\
$10\,000$ & $23\,956$ & $0{,}172$ & $0{,}006$\\
$100\,000$ & $240\,287$ & $0{,}180$ & $0{,}001$\\
\midrule
\textbf{exakt} & & $\mathbf{0{,}1785}$ & \\
\bottomrule
\end{tabular}}
\end{center}

\begin{center}
\begin{minipage}{0.9\textwidth}
\textbf{Die Simulation nähert sich dem exakten Wert --- langsam und unregelmäßig. Bei zehntausend Tagen ist sie schlechter als bei tausend.}
\end{minipage}
\end{center}

\begin{center}
\begin{minipage}{0.88\textwidth}
\kursiv{\textbf{Das ist kein Fehler des Programms, sondern die Natur des Verfahrens}: Jede Zeile der Tabelle ist eine Zufallsgröße. Die dritte war zufällig weit weg, die zweite zufällig nahe. Wer nur eine Zeile sieht, kann das nicht wissen --- und das ist die erste Lehre dieses Bandes: Eine simulierte Zahl ohne Angabe ihrer Streuung ist nicht zu gebrauchen.}
\end{minipage}
\end{center}
\end{merksatz}

\section{Die Streuung messen}

\begin{merksatz}
\boxtitel{Fünfzig Wiederholungen}
Die Simulation über tausend Tage, fünfzigmal mit verschiedenen Startwerten:

\begin{center}
{\small
\renewcommand{\arraystretch}{1.5}
\begin{tabular}{@{}P{4.0cm}P{2.6cm}@{}}
\toprule
Mittel der $50$ Läufe & $0{,}1797$\\
Streuung der Einzelläufe & $0{,}0128$\\
Standardfehler des Mittels & $\mathbf{0{,}0018}$\\
\midrule
exakt & $0{,}1785$\\
\bottomrule
\end{tabular}}
\end{center}

\begin{center}
\begin{minipage}{0.9\textwidth}
\textbf{Das Ergebnis lautet nicht $0{,}18$, sondern $0{,}180 \pm 0{,}002$. Erst mit dem zweiten Teil ist es eine Aussage.}
\end{minipage}
\end{center}

\begin{center}
\begin{minipage}{0.88\textwidth}
\kursiv{\textbf{Ein einzelner Lauf von tausend Tagen streut um $0{,}013$} --- das ist ein Zehntel des Werts selbst. Wer zwei Stationsvarianten mit je einem Lauf vergleicht und einen Unterschied von einem Prozentpunkt findet, hat nichts gefunden. Kapitel~4 zeigt, wie man solche Vergleiche richtig anstellt, und Kapitel~5, wie man die Streuung um das Zehnfache drückt, ohne länger zu rechnen.}
\end{minipage}
\end{center}
\end{merksatz}

\begin{definition}
\boxtitelm{Definition 1.1 --- Monte-Carlo-Schätzer}
Sei $\theta = \mathbb{E}[g(X)]$ für eine Zufallsgröße $X$, die man erzeugen kann. Der \Begriff{Monte-Carlo-Schätzer} ist
\[ \hat\theta_n = \frac1n\sum_{i=1}^{n} g(X_i) \]
mit unabhängigen Kopien $X_1,\ldots,X_n$.
\end{definition}

\begin{merksatz}
\boxtitel{Satz 1.2 --- die Wurzelregel}
Ist $\operatorname{Var}[g(X)] = \sigma^{2} < \infty$, so gilt
\[ \mathbb{E}[\hat\theta_n] = \theta, \qquad \operatorname{Var}[\hat\theta_n] = \frac{\sigma^{2}}{n}, \qquad \frac{\hat\theta_n - \theta}{\sigma/\sqrt n} \ \longrightarrow\ \mathcal N(0,1) . \]

\schritt{Beweis}
Erwartungswert und Varianz folgen aus der Linearität und der Unabhängigkeit; die Grenzverteilung ist der zentrale Grenzwertsatz aus Band 4, Kapitel~13. $\blacksquare$

\begin{center}
\begin{minipage}{0.9\textwidth}
\textbf{Der Fehler fällt mit $1/\sqrt n$ --- unabhängig davon, was $X$ ist, wie viele Dimensionen es hat, und wie glatt $g$ ist.}
\end{minipage}
\end{center}

\begin{center}
\begin{minipage}{0.88\textwidth}
\kursiv{\textbf{Beides ist in diesem Satz enthalten --- der Fluch und der Segen der Methode.} Der Fluch: Zehnmal genauer verlangt hundertmal mehr Läufe, immer. Der Segen: Es gibt keine Abhängigkeit von der Dimension. Die Trapezregel aus Band 5 ist in einer Dimension unschlagbar und in zwanzig unbrauchbar; Monte Carlo ist überall gleich mittelmäßig --- und deshalb oberhalb von etwa fünf Dimensionen die einzige Methode, die überhaupt funktioniert. Kapitel~3 rechnet das vor.}
\end{minipage}
\end{center}
\end{merksatz}

\section{Warum man simuliert}

\begin{merksatz}
\boxtitel{Drei Gründe}
\begin{center}
{\small
\renewcommand{\arraystretch}{1.55}
\begin{tabular}{@{}P{2.6cm}P{6.0cm}@{}}
\toprule
\textbf{Keine Formel} & Sobald die Station Dringlichkeitsstufen hat, die Verweildauer von der Auslastung abhängt oder die Ankünfte tagsüber häufiger sind, gibt es keine Erlang-Formel mehr.\\
\textbf{Verteilung statt Mittel} & Die Formel gibt die mittlere Abweisungsrate. Die Simulation gibt auch, wie oft an einem Tag drei Patienten abgewiesen werden.\\
\textbf{Prüfung} & Ein Schätzer, ein Test, ein Verfahren --- ob es tut, was es soll, prüft man an Daten, bei denen man die Wahrheit kennt. Band 26 hat das durchgehend getan.\\
\bottomrule
\end{tabular}}
\end{center}

\begin{center}
\begin{minipage}{0.88\textwidth}
\kursiv{\textbf{Der dritte Grund ist für die Wissenschaft der wichtigste.} Eine Methode, die in der Simulation unter ihren eigenen Annahmen versagt, ist falsch --- das ist ein Beweis. Eine Methode, die dort funktioniert, ist nicht bewiesen; sie hat nur eine Prüfung bestanden. Das ist der Unterschied zwischen Band B und diesem Band, und er wird in Kapitel~13 wieder aufgenommen.}
\end{minipage}
\end{center}
\end{merksatz}

\section{Der Bruchfall}

\begin{warnung}
\boxtitelr{Was die Simulation nicht weiß}
\begin{center}
{\small
\renewcommand{\arraystretch}{1.52}
\begin{tabular}{@{}P{2.6cm}P{6.0cm}@{}}
\toprule
\textbf{Die Annahmen} & Poisson-Ankünfte, logarithmisch normale Dauern, Unabhängigkeit --- die Simulation nimmt sie als gegeben. Ob sie stimmen, sagt sie nicht.\\
\textbf{Der Anfang} & Die Station startet leer. Die ersten Tage sind untypisch, und bei hundert Tagen Laufzeit sind sie ein Zehntel der Daten.\\
\textbf{Die seltenen Fälle} & Will man wissen, wie oft an einem Tag fünf Patienten abgewiesen werden, braucht man Millionen Tage --- oder Kapitel~6.\\
\bottomrule
\end{tabular}}
\end{center}

\schritt{Zur zweiten Zeile}
\begin{center}
\begin{minipage}{0.88\textwidth}
\kursiv{Die erste Zeile der Tabelle in Abschnitt 1.1 lag bei $0{,}160$ --- zu niedrig, und nicht nur durch Zufall: Eine leere Station weist in der ersten Woche niemanden ab. \textbf{Das ist die Einschwingphase, und sie ist dasselbe Problem wie die Mischzeit in Band 25, Kapitel~3.} Man verwirft die ersten Tage; wie viele, ist die Frage, die Kapitel~7 beantwortet.}
\end{minipage}
\end{center}
\end{warnung}

\section{Übungsaufgaben}

\begin{aufgabe}
\boxtitel{Aufgabe 1.1 --- Die Formel nachrechnen}
Berechnen Sie die Abweisungsrate für $a = 11{,}49$ und $c = 12$ mit der rekursiven Form $B_c = \dfrac{a\,B_{c-1}}{c + a\,B_{c-1}}$, $B_0 = 1$.

\loesung
\schritt{Schritt 1 --- Die Rekursion}
\begin{center}
{\small
\renewcommand{\arraystretch}{1.4}
\begin{tabular}{@{}P{1.2cm}P{2.0cm}P{1.2cm}P{2.0cm}@{}}
\toprule
$c$ & $B_c$ & $c$ & $B_c$\\
\midrule
$0$ & $1{,}000$ & $7$ & $0{,}470$\\
$1$ & $0{,}920$ & $8$ & $0{,}403$\\
$2$ & $0{,}841$ & $9$ & $0{,}340$\\
$3$ & $0{,}763$ & $10$ & $0{,}281$\\
$4$ & $0{,}687$ & $11$ & $0{,}227$\\
$5$ & $0{,}612$ & $12$ & $\mathbf{0{,}1785}$\\
$6$ & $0{,}540$ & $13$ & $0{,}136$\\
\bottomrule
\end{tabular}}
\end{center}

\schritt{Schritt 2 --- Das dreizehnte Bett}
\[ 0{,}1785 - 0{,}136 = 0{,}042 \]

\begin{center}
\begin{minipage}{0.9\textwidth}
\textbf{Ein weiteres Bett senkt die Abweisung von $17{,}9$ auf $13{,}6$ Prozent --- bei $876$ Ankünften im Jahr sind das $37$ Patienten weniger, die verlegt werden müssen.}
\end{minipage}
\end{center}

\schritt{Schritt 3 --- Warum die Rekursion}
\begin{center}
\begin{minipage}{0.88\textwidth}
\kursiv{Die geschlossene Formel verlangt $12!$ und $a^{12}$ --- Zahlen, die bei großen Stationen den Zahlenbereich sprengen. \textbf{Die Rekursion vermeidet das und ist nebenbei numerisch stabil} --- Band 5, Kapitel~2. Dass man diese Aufgabe exakt lösen kann, ist die Ausnahme; sie dient als Maßstab für alles, was folgt.}
\end{minipage}
\end{center}
\end{aufgabe}

\begin{aufgabe}
\boxtitel{Aufgabe 1.2 --- Wie viele Läufe?}
Die Abweisungsrate soll auf $\pm 0{,}002$ genau bestimmt werden ($95\,\%$-Intervall). Wie viele Simulationstage braucht man?

\loesung
\schritt{Schritt 1 --- Die Streuung je Ankunft}
Jede Ankunft ist ein Bernoulli-Versuch mit $p \approx 0{,}18$:
\[ \sigma = \sqrt{0{,}18\cdot0{,}82} = 0{,}384 \]

\schritt{Schritt 2 --- Die Wurzelregel}
\[ 1{,}96\cdot\frac{0{,}384}{\sqrt n} \le 0{,}002 \quad\Longrightarrow\quad n \ge \left(\frac{1{,}96\cdot0{,}384}{0{,}002}\right)^{2} = 141\,600 \ \text{Ankünfte} \]

\schritt{Schritt 3 --- In Tagen}
\[ \frac{141\,600}{2{,}4} \approx 59\,000 \ \text{Tage} \approx 160 \ \text{Jahre} \]

\begin{center}
\begin{minipage}{0.9\textwidth}
\textbf{Für drei Stellen Genauigkeit braucht man hundertsechzig simulierte Jahre --- und die Rechnung unterschätzt noch, weil aufeinanderfolgende Ankünfte nicht unabhängig sind.}
\end{minipage}
\end{center}

\begin{center}
\begin{minipage}{0.88\textwidth}
\kursiv{\textbf{Die Unabhängigkeit fehlt, weil die Station ein Gedächtnis hat}: Ist sie heute voll, ist sie es morgen wahrscheinlich auch. Die Abweisungen kommen in Häufungen, und die wirksame Zahl der Beobachtungen ist kleiner als $n$ --- Satz 10.4 aus Band 26. Die Tabelle aus Abschnitt 1.2 zeigt es: $0{,}0128$ Streuung bei $2387$ Ankünften entspricht einer wirksamen Zahl von nur etwa $900$. Kapitel~4 behandelt die Abhilfe.}
\end{minipage}
\end{center}
\end{aufgabe}

\begin{aufgabe}
\boxtitel{Aufgabe 1.3 --- Eine Frage ohne Formel}
Formulieren Sie die Simulation so um, dass Notfälle Vorrang haben: Ein Notfall belegt das Bett eines Patienten, der am ehesten verlegt werden kann. Was ändert sich am Programm, und was lässt sich nun fragen?

\loesung
\schritt{Schritt 1 --- Was sich ändert}
\begin{center}
{\small
\renewcommand{\arraystretch}{1.5}
\begin{tabular}{@{}P{2.6cm}P{6.0cm}@{}}
\toprule
\textbf{Ankunft} & Jeder Patient erhält eine Dringlichkeit, etwa $30$ Prozent Notfälle.\\
\textbf{Belegung} & Ist alles voll und kommt ein Notfall, wird der Patient mit der kürzesten Restdauer verlegt.\\
\textbf{Buchführung} & Zwei Zähler: abgewiesene Notfälle, vorzeitig Verlegte.\\
\bottomrule
\end{tabular}}
\end{center}

\schritt{Schritt 2 --- Was man nun fragen kann}
\begin{center}
{\small
\renewcommand{\arraystretch}{1.5}
\begin{tabular}{@{}P{1.0cm}P{7.4cm}@{}}
\toprule
$1$ & Wie oft wird ein Notfall abgewiesen? (Nur, wenn alle Betten von Notfällen belegt sind.)\\
$2$ & Wie viele Patienten werden vorzeitig verlegt, und wie lange vor ihrer Zeit?\\
$3$ & Senkt ein dreizehntes Bett eher die Abweisungen oder die Verlegungen?\\
\bottomrule
\end{tabular}}
\end{center}

\begin{center}
\begin{minipage}{0.9\textwidth}
\textbf{Für keine dieser Fragen gibt es eine Formel. Die Simulation braucht zehn Zeilen mehr.}
\end{minipage}
\end{center}

\schritt{Schritt 3 --- Was gleich bleibt}
\begin{center}
\begin{minipage}{0.88\textwidth}
\kursiv{\textbf{Die Wurzelregel.} Auch die neue Simulation streut, und die Frage nach Fehlerbalken, Einschwingphase und Abhängigkeit stellt sich unverändert. Die Freiheit, jede Regel einbauen zu können, ist der Grund für die Beliebtheit der Methode --- und die Statistik dahinter ist dieselbe wie beim einfachsten Fall. Das ist der Aufbau dieses Bandes: Teil~I liefert den Zufall, Teil~II die Genauigkeit, Teil~III und IV die Modelle.}
\end{minipage}
\end{center}
\end{aufgabe}

\begin{geschichte}
\boxtitelg{Ein Kartenspiel und ein Name}
\Person{Stanisław Ulam} lag 1946 nach einer Krankheit im Bett und legte Patiencen. Er fragte sich, wie groß die Wahrscheinlichkeit ist, dass eine aufgeht --- und fand, dass es leichter sei, hundert Spiele zu legen und zu zählen, als es auszurechnen.

\begin{center}
\begin{minipage}{0.9\textwidth}
\kursiv{Er erzählte es \Person{John von Neumann}, und beide erkannten, dass dasselbe für die Neutronendiffusion in Los Alamos galt. \textbf{Den Decknamen \enquote{Monte Carlo} schlug \Person{Nicholas Metropolis} vor} --- nach dem Spielkasino, in dem Ulams Onkel sein Geld verspielt hatte. Die erste Rechnung lief 1947 auf dem ENIAC; von Neumann schrieb das Programm auf Lochkarten.}
\end{minipage}
\end{center}

\schritt{Und die Vorgeschichte}
\begin{center}
\begin{minipage}{0.9\textwidth}
\kursiv{\Person{Georges-Louis Leclerc de Buffon} hatte 1777 vorgeschlagen, $\pi$ zu bestimmen, indem man Nadeln auf ein Dielenmuster wirft und zählt, wie viele eine Fuge kreuzen --- das Titelbild dieses Bandes in anderer Gestalt. \textbf{\Person{William Gosset} erzeugte 1908 seine $t$-Verteilung, indem er Körpermaße von Verbrechern auf Karten schrieb, mischte und zog}: die erste statistische Simulation, von Hand, mit dreitausend Karten.}
\end{minipage}
\end{center}

\schritt{Erlang}
\begin{center}
\begin{minipage}{0.9\textwidth}
\kursiv{Die Formel aus Abschnitt 1.1 stammt von \Person{Agner Krarup Erlang} aus dem Jahr 1917 --- für Telefonleitungen, nicht für Betten. \textbf{Dass sie von der Form der Verweildauer unabhängig ist, bewies erst \Person{Boris Sevastyanov} 1957}; Erlang hatte es vermutet. Für Krankenhäuser wurde sie ab den 1960er Jahren benutzt, und sie ist bis heute das einfachste Werkzeug der Bettenplanung.}
\end{minipage}
\end{center}
\end{geschichte}

\begin{forschung}
\boxtitelp{Simulation heute}
\begin{center}
{\small
\renewcommand{\arraystretch}{1.45}
\begin{tabular}{@{}P{3.0cm}P{5.2cm}@{}}
\toprule
\textbf{Digitale Zwillinge} & Ganze Krankenhäuser als Simulation, fortlaufend mit Daten abgeglichen.\\
\textbf{Agentenmodelle} & Jeder Patient, jeder Arzt ein eigenes Programm --- Band 18 für Epidemien.\\
\textbf{Reproduzier\-barkeit} & Simulationen, die ein anderer mit denselben Zahlen wiederholen kann.\\
\bottomrule
\end{tabular}}
\end{center}

\schritt{Zur dritten Zeile}
\begin{center}
\begin{minipage}{0.88\textwidth}
\kursiv{Ein großer Teil veröffentlichter Simulationsstudien lässt sich nicht wiederholen --- weil der Startwert fehlt, die Version des Generators, die Zahl der verworfenen Einschwingtage. \textbf{Seit 2020 verlangen die meisten Fachzeitschriften Code und Startwerte}; in der Medizin ist man noch nicht so weit. Dieser Band legt für jede Zahl beides offen, und Kapitel~13 macht daraus eine Regel.}
\end{minipage}
\end{center}
\end{forschung}

\begin{merksatz}
\boxtitel{Was Sie jetzt können}
\begin{center}
\renewcommand{\arraystretch}{1.5}
\begin{tabular}{@{}P{3.0cm}P{5.6cm}@{}}
\toprule
\textbf{Der Leitfall} & zwölf Betten, $17{,}9$ Prozent Abweisung\\
\textbf{Definition 1.1} & Mittelwert über Zufallsläufe\\
\textbf{Satz 1.2} & Fehler $\sigma/\sqrt n$, dimensionsfrei\\
\textbf{Die Regel} & jede Zahl mit Standardfehler\\
\textbf{Der Bruchfall} & Annahmen, Anfang, seltene Fälle\\
\bottomrule
\end{tabular}
\end{center}

\begin{center}
\begin{minipage}{0.88\textwidth}
\kursiv{Alles beruht darauf, dass der Rechner Zufallszahlen liefert. Das nächste Kapitel zeigt, dass er das gar nicht kann --- und was er stattdessen tut.}
\end{minipage}
\end{center}
\end{merksatz}
